\chapter{Kawasaki Disease}
\index{Kawasaki Disease}
\label{sec:Kawasaki Disease}

\section{Overview}
Kawasaki disease is an acute, self-limited, systemic vasculitis of medium-sized arteries, primarily affecting children under 5 years of age, with a predilection for the coronary arteries, with the highest incidence in Japan (300 per 100,000 children $<$5 years), and approximately 20--25 per 100,000 in North America and Europe.
\section{Causes}
The cause is unknown but suspected to involve an infectious trigger in genetically susceptible individuals, leading to an intense innate and adaptive immune response, with Disease mechanism involving systemic inflammation, endothelial activation, and infiltration of the vessel wall by neutrophils, lymphocytes, and macrophages, with consequent arteritis that can lead to coronary artery aneurysm formation and blood clot formation.     
\section{Risk Factors}

\begin{itemize}[noitemsep]
  \item Genetic predisposition and family history
  \item Advancing age
  \item Lifestyle factors (diet, exercise, smoking, alcohol)
  \item Environmental exposures
  \item Underlying medical conditions
  \item Medication use
\end{itemize}
\section{Signs and Symptoms}

\subsection{Common symptoms}
\begin{itemize}[noitemsep]
  \item \item Pain or discomfort in the affected area    \item Pain or discomfort in the affected area
  \end{itemize}

\subsection{Emergency warning signs}
\begin{itemize}[noitemsep]
  \item Severe or worsening symptoms of kawasaki disease require immediate medical evaluation
\end{itemize}
\section{Progression and Stages}
Clinical To make the diagnosis, doctors look for fever for at least 5 days plus four of five principal features: bilateral conjunctival injection (non-exudative), oral mucosal changes (strawberry tongue, erythematous lips, pharyngeal redness of the skin), polymorphous rash (usually truncal), extremity changes (redness of the skin and swelling of hands/feet in acute phase, periungual desquamation in subacute phase), and cervical lymphadenopathy ($>$1.5 cm, typically unilateral). Treatment typically involves intravenous immunoglobulin (IVIG, 2 g/kg single dose) plus high-dose aspirin in the acute phase, reduced to low-dose antiplatelet dose after defervescence. The condition may progress through different stages of severity over time. Early stages often involve mild or subtle symptoms that may be overlooked. Moderate stages involve more noticeable symptoms affecting daily function. Advanced stages may involve significant impairment and complications.
\section{Complications}
Complications of coronary arteritis include coronary artery aneurysms (occur in 25\% of untreated KD, reduced to $<$5\% with treatment), heart attack, arrhythmias, and sudden death.

\begin{itemize}[noitemsep]
  \item Disease progression leading to worsening symptoms
  \item Secondary infections or injuries
  \item Side effects of treatment
  \item Reduced quality of life
  \item Emotional and psychological impact
\end{itemize}
\section{Diagnosis}
Comprehensive medical history and physical examination Laboratory tests to assess organ function and rule out other conditions Imaging studies to visualize affected structures Specialized testing depending on the affected system Consultation with relevant specialists Monitoring of disease progression over time

\begin{itemize}[noitemsep]
  \item Comprehensive medical history and physical examination
  \item Laboratory tests to assess organ function
  \item Imaging studies to visualize affected structures
  \item Specialized testing depending on the affected system
  \item Consultation with relevant specialists
\end{itemize}
\section{Prevention}
\begin{itemize}[noitemsep]
  \item Maintain a healthy lifestyle with balanced diet and exercise
  \item Avoid known risk factors
  \item Attend regular medical check-ups for early detection
  \item Follow recommended screening guidelines
\end{itemize}
\section{Treatment Options}
Medications to manage symptoms and address underlying mechanisms Lifestyle modifications and self-care measures Physical therapy and rehabilitation Surgical intervention when indicated Regular monitoring and follow-up care Patient education and support services

\begin{itemize}[noitemsep]
  \item Medications to manage symptoms and address underlying mechanisms
  \item Lifestyle modifications and self-care measures
  \item Physical therapy and rehabilitation
  \item Surgical intervention when indicated
  \item Regular monitoring and follow-up care
\end{itemize}
\section{Living with It}
Echocardiography is essential at diagnosis and follow-up.

\begin{itemize}[noitemsep]
  \item Follow your treatment plan as prescribed
  \item Maintain regular follow-up with your healthcare provider
  \item Make necessary lifestyle adjustments
  \item Seek support from family, friends, or support groups
  \item Stay informed about your condition
\end{itemize}
\section{Outlook}
Most people recover well with treatment, with coronary artery aneurysm rates reduced to $<$5\%. The outlook varies depending on the specific condition, severity at diagnosis, and response to treatment. Early intervention generally leads to better outcomes. Many conditions can be effectively managed with appropriate treatment. Regular follow-up care is important for monitoring and treatment adjustment.
